“Ik wil niet wat vrijkomt
tussen onze band komt.”

was wat Svenska schuw zei. “Er is niet meer een “wij””,

dacht daarop de jager.
Ze was geen lastdrager

en wilde afschuiven
de te zure druiven.

De verrezen horden
konden te hoog worden.

Ze voorzag het effect
dat er werd opgewekt

door haar lepe daden.
Maar voor deze kwaden

verantwoordelijkheid
nemen kon deze meid

totaal niet opbrengen.
Het patroon verlengen

van verborgen hoogmoed
deed de jager niet goed.

“Hoezo? Ga je iets doen
buiten normaal fatsoen?”

Ze verloor haar balans.
Svenska miste de kans

om slim te maskeren
dat ze wilde keren
wie-heeft-hier-de-schuldvraag.
Dat deed de narcist graag.

“Ik heb je gewaarschuwd”,
waarna ze verafschuwd

elk negatief gevoel
of elk ander gewoel.

De jager dacht terug.
Toen ging het niet zo stug,

maar bracht zij juist wat goeds met haar mandje vol zoets.

Op haar blote voetjes
nam Svenska de zoetjes

mee in haar rechterhand
over het mulle zand.

De volle mand wiegde
terwijl zij zweefvliegde

tussen dennenbomen
verdiept in dagdromen.

Ze was dronken van lol
en scheuten alcohol.

Een medebelager
van de mensenjager

stond ineens voor haar neus.
Een gigantische reus.

Ze stond dadelijk stil
voor de bittere pil.

Ze zag de bui hangen;
een blik vol verlangen

keek op het meisje neer.
Er hing een bitse sfeer.

Svenska gaf al spoedig
moedig en koelbloedig

de reus uit haar handje
haar gevulde mandje.

Dat liet hij zich smaken.
Zijn machtige kaken

verbrijzelde koekjes.
Kruimels in de hoekjes

verzamelde de reus
gretig en minutieus.

Svenska sloeg slechts gade
hoe hij zich vollaadde.
